\documentclass[10pt, aspectratio=169]{beamer}
\usetheme[lang=en, font=arial, elevatedtitles=false, serifmath=true]{NMBU}

% Use 'lang=...' to set the language: en (English), bm (Bokmål), nn (Nynorsk)
% Keep 'font=arial' as this is the standard for NMBU
% Use 'elevatedtitles=true' to move the titles up
% Use 'serifmath=false' for Arial math (not the usual LaTeX math)

% Math & units
\usepackage{amsmath, amssymb}
\usepackage[per-mode=symbol]{siunitx}

% Better tables
\usepackage{booktabs}

% Title data
% Title page and metadata
\title{FYSXXX Course Title}
\subtitle{Chapter Y - Chapter Title}
\author{Your Name}
%\institute{Norwegian University of Life Sciences}
\date{\today}

\begin{document}

%========================
\begin{frame}[plain, noframenumbering]  %make sure to have plain and noframenumbering here
  \titlepage
\end{frame}

%========================
\section*{Overview}
\begin{frame}{How to use this style}{And a few short examples}
\begin{itemize}
  \item This is a \emph{generic} starter deck showing common slide layouts.
  \item Replace placeholder text and shapes with your own content.
  \item Compile with \textbf{XeLaTeX} or \textbf{LuaLaTeX} to use the theme's Arial font settings (in Overleaf you can set this in the \textit{File} \(\to\) \textit{Settings} \(\to\) \textit{Compiler}).
  \item Keep this file in the same folder as \texttt{beamerthemeNMBU.sty}.
\end{itemize}
\end{frame}

%========================
\section{Single-column text}
\begin{frame}{Single-column text}
This slide demonstrates a simple text layout.
Paragraphs are set for readability on large screens.
Use concise sentences and limit the number of bullets per slide to keep attention on key points.

Inline math like \(a^2+b^2=c^2\) and short expressions \(x \in \mathbb{R}\) render clearly.
If you prefer the math typesetting to be in Arial, change the \texttt{documentclass} option '\texttt{serifmath=false}'.


Hyperlinks or URLs: \url{https://example.org}.
\end{frame}

%========================
\section{Frame titles}
\begin{frame}[fragile=singleslide]{Frame titles}{And subtitles}

\begin{itemize}
    \item a \texttt{frame title} can be added in two manners: 
    \begin{itemize} 
       \item directly to the frame after the \verb|\begin{frame}| as \verb|\begin{frame}{Frame title}|.
       \item using the \verb|\frametitle{Frame title}| option after \verb|\begin{frame}|.
    \end{itemize}
    \item a \texttt{frame subtitle} can be added in two manners: 
    \begin{itemize} 
       \item directly to the frame after the \verb|\begin{frame}{Frame title}| as \verb|\begin{frame}{Frame title}{Frame subtitle}|.
       \item using the \verb|\framesubtitle{Frame subtitle}| option after \verb|\begin{frame}| and usually after the frame title.
    \end{itemize}
\end{itemize}

\end{frame}

%========================
\section{Two-column layout}
\begin{frame}{Two-column layout}
\begin{columns}[T,onlytextwidth]
  \column{0.52\textwidth}
  \textbf{Left column}
  \begin{itemize}
    \item Use columns for comparisons, pros/cons, or text vs. figures.
    \item Keep lists short; split across slides rather than shrinking fonts.
    \item Use blocks for emphasis (see next slides).
  \end{itemize}

  \column{0.44\textwidth}
  \textbf{Right column}
  \begin{itemize}
    \item Why not more text?
    \item I'm sure by now you are not reading this
  \end{itemize}
\end{columns}
\end{frame}

%========================
\section{Text + figure side-by-side}
\begin{frame}{Text with figure}
\begin{columns}[T,onlytextwidth]
  \column{0.58\textwidth}
  \textbf{Guidelines}
  \begin{enumerate}
    \item Introduce the figure in one or two sentences.
    \item Reference key features or trends your audience should notice.
    \item Conclude with the implication or takeaway.
  \end{enumerate}

  Example sentence with units using \texttt{siunitx}: ``The device length was \SI{120}{mm} and the sampling rate was \SI{10}{kHz}.''

  \column{0.38\textwidth}
  \begin{center}
    % Replace with your image
    % \includegraphics[width=\linewidth]{path/to/image}
    \begin{tikzpicture}
      \draw[NMBUAccent3,dashed,fill=NMBUAccent3!12] (0,0) rectangle (5,4);
      \node at (2.5,2) {Image Placeholder};
    \end{tikzpicture}
  \end{center}
\end{columns}
\end{frame}

%========================
\section{Blocks and emphasis}
\begin{frame}{Blocks and emphasis}
\begin{block}{Key point}
Use a standard block for definitions, main ideas, or neutral emphasis.
\end{block}

\begin{alertblock}{Caution}
Use an alert block for warnings, limitations, or important exceptions.
\end{alertblock}

\begin{exampleblock}{Example}
Use an example block for worked examples, case studies, or brief demonstrations.
\end{exampleblock}
\end{frame}

%========================
\section{Math and units}
\begin{frame}{Equations and units}
\begin{align}
  f(x) &= ax^2 + bx + c, \\
  x_{1,2} &= \frac{-b \pm \sqrt{b^2 - 4ac}}{2a}. \label{eq:quadratic}
\end{align}
Use \texttt{siunitx} for consistent numbers and units: \SI{3.1416}{rad}, \SI{9.81}{m/s^2}, \SI{2.4}{GHz}.
\end{frame}

%========================
\section{Tables}
\begin{frame}{Compact table}
\begin{table}
  \centering
  \small
  \begin{tabular}{@{}lcc@{}}
    \toprule
    Item & Quantity & Note \\
    \midrule
    Alpha & \num{42} & Example entry \\
    Beta & \num{7} & Another entry \\
    Gamma & \num{13} & One more \\
    \bottomrule
  \end{tabular}
  \caption{Compact table using \texttt{booktabs}.}
\end{table}
\end{frame}

%========================
\section{Lists}
\begin{frame}{Lists}
\begin{columns}[T,onlytextwidth]
  \column{0.48\textwidth}
  \textbf{Bulleted list}
  \begin{itemize}
    \item First point
    \item Second point
    \item Third point
  \end{itemize}

  \column{0.48\textwidth}
  \textbf{Numbered list}
  \begin{enumerate}
    \item Step one
    \item Step two
    \item Step three
  \end{enumerate}
\end{columns}
\end{frame}

%========================
\section{Figure slide}
\begin{frame}{Full-width figure}
\begin{figure}
  \centering
  % \includegraphics[width=.9\linewidth]{path/to/figure}
  \setlength{\fboxsep}{0pt}%
  \fbox{\rule{0pt}{5cm}\rule{.9\linewidth}{0pt}}
  \caption{Insert a high-resolution figure with a concise caption.}
\end{figure}
\end{frame}

%========================
\section{Quotations}
\begin{frame}{Quotation}
\begin{quote}
“Concise quotations can provide context or motivation. Keep them short and readable.”
\end{quote}
\end{frame}

%========================
\section{Wrap-up}
\begin{frame}{Takeaways}
\begin{itemize}
  \item One message per slide; visuals over text walls.
  \item Use blocks for emphasis; keep tables compact.
  \item Prefer multiple simple slides to one overloaded slide.
\end{itemize}
\end{frame}

%========================
\begin{frame}
  \centering\LARGE Questions?
\end{frame}

%========================

\section{Backup}
\begin{frame}{Extra slides}
Add additional material here for Q\&A or reference.
\end{frame}

\end{document}
